\documentclass[a4paper,11pt]{letter}
\usepackage[utf8]{inputenc}
\usepackage[T1]{fontenc}
\usepackage{lmodern}
\usepackage[english]{babel}
\usepackage[colorlinks=true, urlcolor=blue]{hyperref}
\usepackage{geometry}
\geometry{left=1in, right=1in, top=1in, bottom=1in}

\name{Baha Khammari}
\address{D\"usseldorf, Germany \\
    \href{mailto:b.e.khammari@gmail.com}{b.e.khammari@gmail.com} \\
    +49\,1520\,9016956}

\begin{document}

\begin{letter}{%
    Bonn Graduate School of Economics (BGSE) \\
    Department of Economics \\
    University of Bonn \\
    Adenauerallee 24--42, 53113 Bonn
}

\opening{Dear Members of the Selection Committee,}

I am writing to apply for the doctoral programme at the Bonn Graduate School of Economics. I am completing my M.Sc.\ in Economics at the University of Cologne (expected September 2026, current GPA: 1.6/1.0) and will be available from October 2026. My research interests are in applied microeconomics, causal inference, and education/labour economics, with a methodological focus on difference-in-differences estimation and heterogeneous treatment effects.

My doctoral project asks which higher education access mechanism, whether scholarships, loans, or affirmative-action quotas, is most effective at reducing labour market inequality, and for whom. I use Brazil as an empirical laboratory because it implemented all three at national scale within a fifteen-year window, but the research questions and methods are designed to inform the same policy choices in Germany, the EU, and other countries. The three mechanisms I study map directly onto instruments used here, namely BAf\"oG, the Deutschlandstipendium, and the ongoing debate about social quotas in university admissions.

My master's thesis (grade: 1.0/1.0, supervised by Anna Person, M.Sc.) provides the first causal analysis of Brazil's ProUni scholarship programme on formal-sector wages at the microregion level. I constructed a panel of 558 regions over 2005--2019 from administrative microdata ($\sim$5\,million RAIS observations) and applied the Callaway \& Sant'Anna (2021) heterogeneity-robust staggered DiD estimator with doubly-robust inference. The key finding, a non-monotonic dose--response with significant wage gains in low-exposure regions that vanish at higher intensities, was entirely masked by conventional two-way fixed effects. The project required managing large-scale administrative datasets via SQL/BigQuery and implementing recent econometric methods in R.

The BGSE is my first choice because its research environment matches my project on several dimensions. Prof.\ Teodora Boneva's work on how beliefs and perceived returns shape higher education decisions, and her finding that the perceived consumption value of education explains much of the socioeconomic gap in enrolment intentions, connects directly to the demand-side mechanisms my project studies from the supply side. Combining her perspective on why students choose (or forgo) higher education with my analysis of how different access mechanisms affect labour market outcomes would open a productive line of work. Her ERC Starting Grant on gender inequality in labour markets also intersects with the gender asymmetry I document in my thesis (persistent gaps despite majority-female scholarship uptake).

Prof.\ Amelie Schiprowski's research on labour market policy evaluation using administrative data and Prof.\ Hans-Martin von Gaudecker's applied microeconometric work on wage inequality provide methodological and substantive overlap with my own approach. More broadly, the ECONtribute Cluster of Excellence, with its focus on markets and public policy, and the BGSE's partnership with IZA create an environment where the labour economics and policy evaluation dimensions of my project would both be well supported.

Before the thesis, I worked for two years at PwC Germany (Assurance Solutions), where I built variance analyses on large financial datasets. An Erasmus+ semester at RSM Erasmus University (Rotterdam) added coursework in asset pricing and derivatives valuation.

My academic CV, a research proposal, and a thesis abstract are enclosed. References are available on request.

Thank you for considering my application.

\closing{Yours sincerely,}

\end{letter}

\end{document}
